\section{So sánh và đánh giá hiệu suất các mô hình}
\label{sec:model_comparison_detailed}

Trong khuôn khổ của nghiên cứu này, bốn mô hình hồi quy đã được nhóm em xây dựng và đánh giá nhằm mục tiêu dự đoán sản lượng cây trồng (\texttt{Crop\_Yield\_MT\_per\_HA}) dựa trên các yếu tố về thời tiết, nông nghiệp, và các thông tin liên quan khác. Các mô hình bao gồm:
\begin{enumerate}
    \item Hồi quy Tuyến tính (Linear Regression)
    \item Cây Quyết định Regressor (Decision Tree Regressor) với các tham số mặc định
    \item Mạng Nơ-ron Nhân tạo (Artificial Neural Network - ANN) sử dụng MLPRegressor với cấu hình ban đầu
    \item K-nearest Neightbors (K-Nearest Neighbors - KNN) Regressor với K=5
\end{enumerate}
Tất cả các mô hình đều sử dụng cùng một quy trình tiền xử lý dữ liệu cơ bản (xử lý giá trị thiếu, mã hóa One-Hot cho biến phân loại). Đối với mô hình ANN và KNN, các đặc trưng đầu vào đã được chuẩn hóa bằng \texttt{StandardScaler} do tính nhạy cảm của các thuật toán này với thang đo của dữ liệu.

Phần này nhóm em sẽ trình bày so sánh chi tiết hiệu suất của các mô hình dựa trên các chỉ số đánh giá hồi quy phổ biến: R-squared (R²), Mean Absolute Error (MAE), Mean Squared Error (MSE), và Root Mean Squared Error (RMSE), được tính toán trên tập dữ liệu kiểm tra.

\subsection{Tổng hợp dết quả đánh giá hiệu suất}
\label{subsec:summary_evaluation_metrics_detailed}

Bảng \ref{tab:model_comparison_metrics_detailed} tóm tắt các chỉ số hiệu suất chính của từng mô hình. Giá trị R² phản ánh tỷ lệ phần trăm phương sai của biến mục tiêu mà mô hình có thể giải thích được. Các chỉ số MAE, MSE, và RMSE đo lường mức độ sai số trung bình của các dự đoán; giá trị càng thấp càng tốt.

\begin{table}[H] % Sử dụng [H] từ gói float để ép bảng ở đây (cẩn thận với bố cục)
    \centering
    \caption{Bảng so sánh chi tiết các chỉ số đánh giá hiệu suất của các mô hình hồi quy trên tập kiểm tra. Các giá trị lỗi (MAE, RMSE) được hiểu theo đơn vị của biến mục tiêu (ví dụ: tấn/ha).}
    \label{tab:model_comparison_metrics_detailed}
    \resizebox{\textwidth}{!}{% Bảng tự động co giãn theo chiều rộng trang
    \begin{tabular}{l c c c c}
        \toprule
        \textbf{Tên Mô hình} & \textbf{R-squared (R²)} & \textbf{MAE} & \textbf{MSE} & \textbf{RMSE} \\
        \midrule
        Hồi quy Tuyến tính & $0.5596$ & $0.5466$ & $0.4648$ & $0.6818$ \\
        Cây Quyết định Regressor (tham số gốc) & $0.2679$ & $0.6705$ & $0.7727$ & $0.8791$ \\
        Mạng Nơ-ron Nhân tạo (ANN - MLPRegressor)$^*$ & $0.5218$ & $0.5614$ & $0.5047$ & $0.7104$ \\
        K-nearest Neightbors (KNN, K=5)$^*$ & $0.1284$ & $0.7771$ & $0.9200$ & $0.9592$ \\
        \bottomrule
    \end{tabular}%
    }
    \footnotesize{$^*$ Các đặc trưng đầu vào cho mô hình ANN và KNN đã được chuẩn hóa bằng StandardScaler.}
\end{table}

\subsection{Phân tích chi tiết và thảo luận Kết quả}
\label{subsec:comparative_analysis_detailed}

Từ Bảng \ref{tab:model_comparison_metrics_detailed}, chúng tôi (trong báo cáo bạn sẽ dùng "nhóm em") tiến hành phân tích sâu hơn về hiệu suất của từng mô hình:

\subsubsection{Hồi quy Tuyến tính (Linear Regression)}
Mô hình Hồi quy Tuyến tính đạt được giá trị $R^2 = 0.5596$, cao nhất trong số bốn mô hình được thử nghiệm. Điều này có nghĩa là khoảng $55.96\%$ sự biến thiên trong sản lượng cây trồng có thể được giải thích bởi các đặc trưng đầu vào thông qua một mối quan hệ tuyến tính. Các chỉ số lỗi MAE ($0.5466$) và RMSE ($0.6818$) cũng là thấp nhất, cho thấy độ chính xác dự đoán tương đối tốt so với các mô hình khác trong thử nghiệm này.
\begin{itemize}
    \item \textbf{Ưu điểm:} Đơn giản, dễ triển khai, kết quả dễ diễn giải thông qua các hệ số hồi quy (như đã trình bày ở mục \ref{subsubsec:lr_coefficients}).
    \item \textbf{Hạn chế:} Giả định mối quan hệ tuyến tính giữa các đặc trưng và biến mục tiêu, có thể không nắm bắt được các tương tác phức tạp hoặc các mối quan hệ phi tuyến.
    \item \textbf{Đồ thị trực quan (Mục \ref{subsubsec:lr_visualization}):} Biểu đồ thực tế và dự đoán (Hình \ref{fig:lr_actual_vs_predicted}) cho thấy các điểm bám khá tốt quanh đường chéo, và phân phối phần dư (Hình \ref{fig:lr_residuals_distribution}) tương đối cân đối quanh 0.
\end{itemize}

\subsubsection{Mạng Nơ-ron nhân tạo (ANN - MLPRegressor)}
Mô hình ANN, với kiến trúc ban đầu gồm hai lớp ẩn (100 và 50 nơ-ron) và cơ chế dừng sớm (dừng sau 22 vòng lặp), đạt $R^2 = 0.5218$. Kết quả này rất đáng khích lệ, chỉ thấp hơn một chút so với Hồi quy Tuyến tính. MAE ($0.5614$) và RMSE ($0.7104$) cũng ở mức cạnh tranh.
\begin{itemize}
    \item \textbf{Ưu điểm:} Có khả năng học các mối quan hệ phi tuyến phức tạp mà không cần định nghĩa tường minh các đặc trưng tương tác hay đa thức. Đường cong hàm mất mát (Hình \ref{fig:ann_loss_curve}) cho thấy mô hình đã học hiệu quả.
    \item \textbf{Hạn chế:} Là một mô hình "hộp đen", khó diễn giải trực tiếp các trọng số. Hiệu suất rất nhạy cảm với kiến trúc mạng và các siêu tham số (tốc độ học, số epochs, hàm kích hoạt, thuật toán tối ưu, etc.), đòi hỏi quá trình tinh chỉnh kỹ lưỡng. Yêu cầu chuẩn hóa đặc trưng đầu vào.
    \item \textbf{Đồ thị trực quan (Mục \ref{subsubsec:ann_visualization}):} Các biểu đồ (Hình \ref{fig:ann_actual_vs_predicted} và \ref{fig:ann_residuals_distribution}) cũng cho thấy hiệu suất tốt, tương tự Hồi quy Tuyến tính.
\end{itemize}

\subsubsection{Cây quyết định Regressor (Decision Tree Regressor)}
Mô hình cây Quyết định với các tham số mặc định cho kết quả $R^2 = 0.2679$, thấp hơn đáng kể so với hai mô hình trên. Các chỉ số lỗi MAE ($0.6705$) và RMSE ($0.8791$) cũng cao hơn.
\begin{itemize}
    \item \textbf{Ưu điểm:} Dễ hiểu, dễ diễn giải thông qua cấu trúc cây (như Hình \ref{fig:dt_tree_visualization} đã minh họa một phần). Không yêu cầu chuẩn hóa đặc trưng. Có thể nắm bắt các mối quan hệ phi tuyến.
    \item \textbf{Hạn chế:} Một cây quyết định đơn lẻ (không giới hạn độ sâu hoặc không được cắt tỉa) rất dễ bị học quá khớp (overfitting) với dữ liệu huấn luyện, dẫn đến hiệu suất kém trên tập kiểm tra. Độ ổn định không cao, một thay đổi nhỏ trong dữ liệu có thể dẫn đến một cấu trúc cây rất khác.
    \item \textbf{Đồ thị trực quan (Mục \ref{subsubsec:dt_visualization}):} Biểu đồ thực tế và dự đoán (Hình \ref{fig:dt_actual_vs_predicted}) cho thấy sự phân tán lớn hơn và các "bậc thang" trong dự đoán, đặc trưng của cây quyết định.
\end{itemize}
Mặc dù độ quan trọng của đặc trưng (Mục \ref{subsubsec:dt_feature_importances}) cung cấp thông tin hữu ích, hiệu suất tổng thể cho thấy cần phải tinh chỉnh hoặc sử dụng các kỹ thuật ensemble dựa trên cây.

\subsubsection{K-nearest Neighbors (KNN Regressor)}
Mô hình KNN với $K=5$ và trọng số `distance` cho kết quả thấp nhất trong các thử nghiệm, với $R^2 = 0.1284$. MAE ($0.7771$) và RMSE ($0.9592$) cũng là cao nhất.
\begin{itemize}
    \item \textbf{Ưu điểm:} Đơn giản về mặt khái niệm, không có giả định mạnh về phân phối dữ liệu. Có thể hoạt động tốt với các mối quan hệ phi tuyến phức tạp cục bộ.
    \item \textbf{Hạn chế:} Hiệu suất rất phụ thuộc vào việc chọn giá trị $K$ và metric đo khoảng cách. Nhạy cảm với "lời nguyền chiều dữ liệu" (curse of dimensionality), đặc biệt khi có nhiều đặc trưng không liên quan hoặc dư thừa. Yêu cầu chuẩn hóa đặc trưng đầu vào. Chi phí tính toán dự đoán có thể cao với tập dữ liệu lớn do phải tính toán khoảng cách đến tất cả các điểm trong tập huấn luyện.
    \item \textbf{Đồ thị trực quan (Mục \ref{subsubsec:knn_visualization}):} Các biểu đồ (Hình \ref{fig:knn_actual_vs_predicted} và \ref{fig:knn_residuals_distribution}) phản ánh hiệu suất chưa cao của mô hình này.
\end{itemize}

\subsection{Thảo luận chung và Kết luận sơ bộ}
\label{subsec:overall_discussion_preliminary_conclusion}

Dựa trên các kết quả đánh giá ban đầu này, nhóm em kết luận:
\begin{enumerate}
    \item \textbf{Hồi quy Tuyến tính} đã thiết lập một baseline mạnh mẽ. Điều này cho thấy một phần đáng kể của mối quan hệ giữa các yếu tố đầu vào và sản lượng có thể là tuyến tính hoặc gần tuyến tính.
    \item \textbf{Mạng Nơ-ron Nhân tạo (ANN)} cho thấy tiềm năng lớn, đạt hiệu suất gần bằng Hồi quy Tuyến tính ngay cả với cấu hình ban đầu. Việc tinh chỉnh sâu hơn có thể giúp ANN vượt qua Hồi quy Tuyến tính, đặc biệt nếu có các mối quan hệ phi tuyến phức tạp mà mô hình tuyến tính bỏ lỡ.
    \item \textbf{Cây Quyết định và KNN}, với các thiết lập mặc định hoặc ban đầu, chưa cho thấy hiệu quả cao bằng hai mô hình trên. Cả hai mô hình này đều đòi hỏi sự tinh chỉnh siêu tham số cẩn thận (ví dụ: độ sâu của cây, số lượng mẫu tối thiểu để chia nhánh cho Cây Quyết định; giá trị $K$ cho KNN) để có thể đạt được hiệu suất tốt nhất.
\end{enumerate}

Kết quả này nhấn mạnh tầm quan trọng của việc lựa chọn mô hình phù hợp với bản chất của dữ liệu cũng như sự cần thiết của việc tinh chỉnh siêu tham số. Đối với các mô hình phức tạp như ANN, Cây quyết định, và KNN, việc sử dụng các giá trị tham số mặc định thường không mang lại kết quả tối ưu.

